\clearpage
\VerdierPhysicalPageAnchor{0233}%
\begingroup\small\setlength{\parskip}{0pt}\setlength{\emergencystretch}{1.5em}
\setlength{\abovedisplayskip}{4pt plus 2pt minus 2pt}%
\setlength{\belowdisplayskip}{4pt plus 2pt minus 2pt}%

\VerdierSemanticAnchor{ch4d-4-1-8}
\noindent\textbf{4.1.8.}\quad
When $X$ is a simple complex, the spectral object with values in
$K(\mathcal A)$ defined by the complex $X$ filtered by its first degree
(4.1.7) is precisely the first spectral object with values in
$K(\mathcal A)$ associated with the complex $X$.

\VerdierSemanticAnchor{ch4d-4-1-8-degree-support}
Likewise, let $X$ be an $n$-fold complex and let $m$ be an integer,
$1\leq m\leq n$. Suppose that the component $X^\sigma$ of $X$ is zero
whenever
\[
\sigma-\sigma_m e_m\neq0.
\]
Here $\sigma_m$ denotes the $m$-th coordinate.
One then verifies that the spectral object defined by the complex $X$
filtered by its $m$-th degree is equal to the first canonical spectral
object of the simple complex $\int_m X$.

\VerdierSemanticAnchor{ch4d-4-1-8-derived-image}
Finally, when the category $\mathcal A$ is abelian, the canonical functor
\[
K(\mathcal A)\longrightarrow D(\mathcal A)
\]
takes the spectral object with values in $K(\mathcal A)$ defined by an
$n$-fold complex $X$ filtered by its $m$-th degree to a spectral object
with values in $D(\mathcal A)$. The spectral object thus obtained is called
\emph{the spectral object, with values in $D(\mathcal A)$, defined by $X$,
filtered by its $m$-th degree}. When no confusion can result, the phrase
``with values in $D(\mathcal A)$'' is omitted.

\VerdierSemanticAnchor{ch4d-4-1-9}
\medskip\noindent\textbf{4.1.9.}\quad
Let $X^{\bullet\bullet}$ be a double complex in $\mathcal A$.

\VerdierSemanticAnchor{ch4d-4-1-9-row-complexes}
For every integer $q$, denote by $X^{q,\bullet}$ (respectively,
$X^{\bullet,q}$) the simple complex
\[
\begin{aligned}
X^{q,\bullet}
&=\cdots\xrightarrow{d_2^{q,i-1}}X^{q,i}
  \xrightarrow{d_2^{q,i}}X^{q,i+1}
  \xrightarrow{d_2^{q,i+1}}\cdots,\\[3pt]
(\text{resp.}\quad X^{\bullet,q}
&=\cdots\longrightarrow X^{i,q}
  \xrightarrow{(-1)^q d_1^{i,q}}X^{i+1,q}
  \xrightarrow{(-1)^q d_1^{i+1,q}}X^{i+2,q}
  \longrightarrow\cdots\quad),
\end{aligned}
\]
whose component of degree $i$ is $X^{q,i}$ (respectively, $X^{i,q}$).

\VerdierSemanticAnchor{ch4d-4-1-9-direction-morphisms}
The differential in the first (respectively, second) direction of
$X^{\bullet\bullet}$ induces morphisms (in $K(\mathcal A)$):
\[
\begin{aligned}
d_1^{q,\bullet}&\colon X^{q,\bullet}\longrightarrow X^{q+1,\bullet}\\
(\text{resp.}\quad
d_2^{\bullet,q}&\colon X^{\bullet,q}\longrightarrow X^{\bullet,q+1}\quad),
\end{aligned}
\]

\VerdierSemanticAnchor{ch4d-4-1-9-components}
whose components are
\[
\begin{aligned}
(d_1^{q,\bullet})^i&=d_1^{q,i}\\
(\text{resp.}\quad(d_2^{\bullet,q})^i&=(-1)^i d_2^{i,q}\quad).
\end{aligned}
\]
\endgroup
